% Reproduction Steps
\needspace{10cm}
\section{Reproduction Steps}

% Provide a deterministic sequence to trigger the defect.
% Each step should be concrete, observable, and independently verifiable.

\subsection{Preconditions}

% State what must be true before starting the reproduction sequence.

\begin{itemize}
    \item Precondition 1 (e.g., user is authenticated with role X)
    \item Precondition 2 (e.g., database seeded with test fixtures)
    \item Precondition 3 (e.g., feature flag Y is enabled)
\end{itemize}

\subsection{Steps to Reproduce}

\begin{enumerate}
    \item First action — describe exactly what to do
    \item Second action — include specific inputs or data values
    \item Third action — note any timing or ordering constraints
    \item Observe the defect at this step
\end{enumerate}

\subsection{Reproduction Rate}

\begin{wyrdtab}{l l}
\textbf{Frequency}       & Always / Intermittent / Rare \\
\textbf{Attempts}        & X out of Y attempts \\
\textbf{Conditions}      & Any special timing, load, or data conditions \\
\end{wyrdtab}

% --- Repeat per bug ---
% When batching, duplicate the Steps to Reproduce subsection per defect:
%
% \subsection{Bug 2: Steps to Reproduce}
% \begin{enumerate}
%     \item ...
% \end{enumerate}
